% !TEX root = ../main.tex
% !TEX program = xelatex
\section{问题分析}

四个问题围绕同一微网系统展开，核心都是在满足小区负载的前提下，协调光伏、外网购电和储能设备的使用。
随着问题推进，决策条件由“供需和价格已知”逐步扩展到“供需未知、预报更新以及电价波动”，因此建模过程也由确定性调度逐步转向预测与滚动优化。

\subsection{问题一的分析}

问题一属于\textbf{确定性优化问题}。题目已经给出一天内的负载、光伏预测和电价，因此本问不需要预测未来，重点是确定各时段何时购电、何时充电以及何时放电。

本问的主要难点是供需与价格存在明显的时间错位。白天可能出现光伏富余，而晚间仍有较大用电需求和较高电价。
若只根据当前电价独立决策，可能过早充满电池，无法继续接收后续光伏；也可能过早放电，导致真正高价时段缺少可用储能。
因此，需要同时考虑全天各时段，而不能逐时独立安排。

此外，题目要求0:00和24:00的储电量相同，从而避免通过消耗原有库存人为降低当天费用。
基于这些特点，本文将全天144个时段统一纳入\textbf{全时段电量流线性规划模型}，
在储能容量、充放电功率和效率等约束下，协调外网购电、光伏利用和储能转移，得到全天购电与充放电方案。

\subsection{问题二的分析}

问题二属于\textbf{预测与优化相结合的决策问题}。
与问题一不同，每天0:00制定计划时，当天真实负载和光伏尚未发生，因此不能直接使用附件中的未来实际数据，而必须根据此前已经发生的历史信息进行估计。

负载和光伏具有不同的变化特点：负载存在较明显的星期重复，而光伏更多受到近期发电水平变化的影响，因此二者分别预测，再合成为未来的净需求。

本问另一个关键困难是购电风险具有明显的不对称性。
普通电订多后即使没有使用，也需要付费；若订少导致供电不足，则需要按正常电价的5倍进行紧急购电。
储能虽然能够缓冲部分预测误差，但仍受到最低储电量和最大放电功率限制，因此不能仅靠电池解决全部偏差。

因此，本文采用\textbf{周期预测—近期加权准备—跨日线性规划}的思路：
先预测负载和光伏，再根据近期曾出现的净需求预测误差设置适量安全准备，最后利用线性规划安排普通购电和储能。
日末储电量直接传递到下一天，使各日运行连续衔接，而不是每天重新开始。

\subsection{问题三的分析}
问题三属于\textbf{动态信息下的滚动优化问题}。相比问题二，系统在6:00、12:00和18:00可以得到新的光伏预报，并允许调整尚未执行的购电计划。
问题的重点因此不只是“新预报是否更准确”，而是“利用新预报修改订单是否真正更经济”。

新的光伏预报可以修正零点时的判断，但改单本身需要支付额外费用。
如果电池仍有足够的调节空间，部分预测变化也可以直接由储能吸收。
因此，收到新预报并不意味着一定要修改订单，还需权衡当前储电量、调整费用和不调整可能产生的紧急购电风险。

负载也会随当天运行逐渐明朗。星期规律可以给出零点时的参考，
但当天用电可能持续偏高或偏低，较早的历史也可能跟不上近期变化。
因此，零点预测应兼顾星期规律与近期性；日内修正则要判断已发生的偏差能否延续，避免把短时波动放大到后续时段。

附件3给出的光伏预报为整点数据，而调度间隔为10分钟，因此需先将预报转换到统一时间尺度。
同时，官方预报能够反映当天信息，近期历史光伏则可以提供稳定参考，因此本文采用\textbf{融合预测与日内校正滚动规划}。
每次决策更新负载判断、融合光伏预测和误差准备，再从当前真实库存优化尚未执行的部分。

此外，在评价日内预报价值时，需要区分“允许日内观察实际运行并修改订单”的作用和“新增光伏预报本身”的作用，避免把两类信息带来的收益混为一谈。

\subsection{问题四的分析}

问题四是在问题二、三的基础上进一步加入波动电价的不确定性。此时未来价格在决策时也尚未发生，因此不能简单沿用原来的购电方案，再用实际电价重新计算费用。因为电价变化本身会改变何时适合购电、何时应该充电以及储能应保留到哪个时段。

本问仍需分别对应问题二和问题三的两种决策权限。

对于问题四-2，每天只能在0:00确定普通购电计划，无法利用当天已经发生的价格进行修正，因此需要从历史价格、星期位置、季节特征和预测净负荷中判断未来价格。
考虑到这些因素与价格之间可能存在非线性关系，采用\textbf{提升树预测与零点规划}。

对于问题四-3，6:00、12:00和18:00可以利用当天已经发生的价格更新后续判断。
由于历史价格特征之间相关性较强，同时需要保持日内多次修正的稳定性和可解释性，
因此采用\textbf{岭回归预测与日内滚动规划}。
供需部分直接承接问题三的近期负载预测、日内校正、光伏融合和误差准备，再加入价格判断。
预测价格只作为购电和储能优化的决策依据，最终费用仍按照真实发生的电价结算。
